% SPDX-License-Identifier: Apache-2.0
% Build with database-de.activities.json (same language).
\documentclass[aspectratio=169,11pt]{beamer}
\usepackage[ngerman]{babel}
\newcommand{\DatabaseFooter}{Datenbank-Aktivitäten}
\newcommand{\DatabaseSlide}[1]{\input{slides-de/#1}}
% SPDX-License-Identifier: Apache-2.0
\usepackage[T1]{fontenc}
\usepackage{lmodern}
\usepackage{texstage}
\usepackage{texquiz}
\usepackage{texactivity}
\usepackage{booktabs}
\usepackage{texstage-coordinates-plane}
\usepackage{texstage-er-decisions}
\usepackage{texstage-relational-tree}
\usepackage{texstage-queryplan-challenge}
\usepackage{texstage-bplus-tree}
\usepackage{texstage-fd-lab}
\usepackage{texstage-normalization-workshop}
\usepackage{texstage-transactions-schedule}
\definecolor{dbink}{HTML}{20364D}
\definecolor{dbteal}{HTML}{087F8C}
\setbeamercolor{normal text}{fg=dbink,bg=white}
\setbeamercolor{structure}{fg=dbteal}
\setbeamercolor{frametitle}{fg=dbink}
\setbeamerfont{frametitle}{size=\Large,series=\bfseries}
\setbeamertemplate{navigation symbols}{}
\setbeamersize{text margin left=9mm,text margin right=9mm}
\setbeamertemplate{footline}{\hfill\scriptsize \DatabaseFooter\quad
  \insertframenumber\hspace{5mm}\vspace{3mm}}

\renewcommand\texactivitylabel{Live-Aktivität -- auf dem Handy}
\renewcommand\texactivityreferencelabel{Referenzlösung}
\quizsetup{time=60,points=none,leaderboard=false,shuffle=false,
  handout-heading={Live-Quiz},
  handout-text={Diese Frage wird live in der Veranstaltung gestellt.},
  results-text={Die Ergebnisse werden live in der Veranstaltung gezeigt.},
  true-label={Wahr},false-label={Falsch},
  number-hint={Geben Sie eine Zahl auf dem Handy ein.},
  text-hint={Geben Sie Ihre Antwort auf dem Handy ein.},
  multiple-hint={Mehrere Antworten können richtig sein.},
  poll-hint={Wählen Sie Ihre Antwort auf dem Handy.},
  order-hint={Bringen Sie die Elemente in die richtige Reihenfolge.},
  sketch-hint={Zeichnen Sie Ihre Antwort auf dem Handy.},
  answer-label={Antwort:},accepted-label={Akzeptierte Antworten:},
  answer-title={Antwort},answer-heading={Antwort},
  answer-text={Die Antwort wird in der Veranstaltung gezeigt.},
  leaderboard-title={Rangliste},podium-title={Podium},
  join-title={Am Datenbanklabor teilnehmen},
  join-text={Öffnen Sie die angegebene Adresse auf dem Handy und geben Sie die PIN ein.}}
\title{Eine Campus-Datenbank, viele Denkwege}
\subtitle{Eingebaute Aktivitäten und alle aktuellen Unterrichtsplugins}
\author{Fiktive Lehrbeispiele}
\date{}
\begin{document}
% SPDX-License-Identifier: Apache-2.0
\DatabaseSlide{00-welcome.tex}
\DatabaseSlide{31-activity-map.tex}
\DatabaseSlide{01-model.tex}
\DatabaseSlide{02-loans.tex}
\DatabaseSlide{03-join.tex}
\DatabaseSlide{32-coordinate-warmup.tex}
\DatabaseSlide{04-primary-key.tex}
\DatabaseSlide{05-constraints.tex}
\DatabaseSlide{06-er-sketch.tex}
\DatabaseSlide{33-er-introduction.tex}
\DatabaseSlide{34-er-registrations.tex}
\DatabaseSlide{07-sql-reminder.tex}
\DatabaseSlide{08-sql-filter.tex}
\DatabaseSlide{09-sql-join.tex}
\DatabaseSlide{10-sql-count.tex}
\DatabaseSlide{11-zero-debrief.tex}
\DatabaseSlide{12-sql-available.tex}
\DatabaseSlide{13-sql-bags.tex}
\DatabaseSlide{14-duplicates.tex}
\DatabaseSlide{15-active-total.tex}
\DatabaseSlide{16-distinct.tex}
\DatabaseSlide{17-query-order.tex}
\DatabaseSlide{35-relational-introduction.tex}
\DatabaseSlide{36-ra-construct.tex}
\DatabaseSlide{37-ra-rewrite.tex}
\DatabaseSlide{38-plan-introduction.tex}
\DatabaseSlide{39-plan-selective.tex}
\DatabaseSlide{40-plan-broad.tex}
\DatabaseSlide{41-plan-library.tex}
\DatabaseSlide{42-plan-memory.tex}
\DatabaseSlide{43-bplus-introduction.tex}
\DatabaseSlide{44-bplus-lookup.tex}
\DatabaseSlide{45-bplus-range.tex}
\DatabaseSlide{46-bplus-insert.tex}
\DatabaseSlide{47-bplus-delete.tex}
\DatabaseSlide{48-bplus-cascade.tex}
\DatabaseSlide{49-bplus-shrink.tex}
\DatabaseSlide{50-fd-introduction.tex}
\DatabaseSlide{51-fd-closure.tex}
\DatabaseSlide{52-fd-practice.tex}
\DatabaseSlide{53-fd-keys-one.tex}
\DatabaseSlide{54-fd-keys-all.tex}
\DatabaseSlide{55-fd-implied.tex}
\DatabaseSlide{56-fd-refuted.tex}
\DatabaseSlide{57-fd-different-sizes.tex}
\DatabaseSlide{58-fd-empty-key.tex}
\DatabaseSlide{59-normalization-introduction.tex}
\DatabaseSlide{60-norm-registration.tex}
\DatabaseSlide{61-norm-classify.tex}
\DatabaseSlide{62-norm-tradeoff.tex}
\DatabaseSlide{18-rdf-reminder.tex}
\DatabaseSlide{19-crowd-facts.tex}
\DatabaseSlide{20-crowd-pairs.tex}
\DatabaseSlide{21-sparql-select.tex}
\DatabaseSlide{22-sparql-ask.tex}
\DatabaseSlide{23-sparql-construct.tex}
\DatabaseSlide{24-turtle-one.tex}
\DatabaseSlide{25-turtle-many.tex}
\DatabaseSlide{29-transaction-introduction.tex}
\DatabaseSlide{63-tx-counter.tex}
\DatabaseSlide{64-tx-independent.tex}
\DatabaseSlide{65-tx-read-read.tex}
\DatabaseSlide{66-tx-write-write.tex}
\DatabaseSlide{67-tx-chain.tex}
\DatabaseSlide{68-tx-blind-write.tex}
\DatabaseSlide{26-word-cloud.tex}
\DatabaseSlide{27-explain.tex}
\DatabaseSlide{28-topic-poll.tex}
\DatabaseSlide{69-feedback-and-points.tex}
\DatabaseSlide{30-takeaways.tex}

\end{document}
